% The four parts of the Day 7 lab. The steps follow Labs/day07/README.md.
%
% Hardware status: not tested on hardware. A value that only the board can
% give is written as "..." and the students measure it. The deck does not
% show a result that the students must predict.

% ------------------------------------------------------------------- Part A
\begin{frame}[fragile]{Part A: setup (35 min)}
  \small
  \renewcommand{\arraystretch}{1.15}
  \begingroup\centering
  \begin{tabular}{@{}L{0.07\textwidth}L{0.70\textwidth}L{0.12\textwidth}@{}}
    \toprule
    \textbf{Step} & \textbf{Work} & \textbf{Time} \\
    \midrule
    1 & Start the board, and connect with SSH & 10 min \\
    2 & Check the system with the script of the course card & 10 min \\
    3 & Test the camera, and look at the photo on the laptop & 10 min \\
    4 & Copy the lab folder to the board & 5 min \\
    \bottomrule
  \end{tabular}\par\endgroup
\begin{lstlisting}[style=shell]
ssh edge@pi-NN.local                      # NN is the number of your group
bash ~/HW-04/check_pi.sh                  # each line: PASS, FAIL, or INFO
rpicam-hello --list-cameras
rpicam-jpeg --output test.jpg --width 640 --height 480
scp -r day07 edge@pi-NN.local:~/edgeai/   # on the laptop, in the folder Labs/
\end{lstlisting}
  \begin{itemize}
    \item \textbf{You record:} the model, the release of the operating
          system, the Python version, the RAM, and the temperature.
    \item Open a second SSH connection and start \code{htop}. Keep it open:
          it shows the load of each of the 4 cores.
    \item Tell the instructor if a line of the check shows \code{FAIL}.
  \end{itemize}
\end{frame}

% ------------------------------------------------------------------- Part B
\begin{frame}[fragile]{Part B: the first inference (35 min)}
  \small
  Get the cat photo of the kit lab, and run the model of the kit lab:
\begin{lstlisting}[style=shell]
wget -P images https://upload.wikimedia.org/wikipedia/commons/3/3a/Cat03.jpg
python pi/classify_image.py images/Cat03.jpg
\end{lstlisting}
\begin{lstlisting}
Model:  models/mobilenet_v2_1.0_224_quant.tflite (3577760 bytes)
Image:  images/Cat03.jpg, threads: 4
Input:  shape [1, 224, 224, 3], type uint8
Output: 1001 values, type uint8
Load: ... ms. First inference: ... ms. Median of the next runs: ... ms.
Task B1: not complete. The list shows the raw output values.
\end{lstlisting}
  \begin{itemize}
    \item The model is MobileNetV2 in \code{uint8}, with 1001 classes.
    \item The script reads the shape and the type of the input from the
          file, and prepares the image for them.
  \end{itemize}
  \begin{exerciseblock}
    \textbf{Prediction.} With one MAC in each clock cycle, this model needs
    125 ms on the board. Which median latency do you expect with 4 threads?
    Write your number in the report. Then run the script.
  \end{exerciseblock}
  \source{Photo: Fir0002, Wikimedia Commons.}
\end{frame}

\begin{frame}{Part B: Task B1 and the camera}
  \small
  \renewcommand{\arraystretch}{1.2}
  \begingroup\centering
  \begin{tabular}{@{}L{0.07\textwidth}L{0.70\textwidth}L{0.12\textwidth}@{}}
    \toprule
    \textbf{Step} & \textbf{Work} & \textbf{Time} \\
    \midrule
    1, 2 & Get the photo, and run the model & 15 min \\
    3 & Task B1: from the raw outputs to probabilities & 10 min \\
    4 & Classify photos of the camera, and compare 1 thread with 4 threads
      & 10 min \\
    \bottomrule
  \end{tabular}\par\endgroup
  \vskip 0.4em
  \begin{itemize}
    \item \textbf{Task B1} is the function \code{dequantize\_and\_softmax}
          in \code{pi/classify\_image.py}. The mapping of Day 4 gives the
          real value: $r = S\,(q - Z)$. Then apply the softmax.
    \item The check: the script prints \code{Task B1: complete}, and a cat
          class is first. On the work computer, the first class has 39
          percent.
    \item The camera:
          \code{python pi/classify\_camera.py --count 3 --interval 3}.
          Point it at a keyboard, a mouse, a cup, or a bottle.
    \item \textbf{You record:} the five classes of the cat photo, and for
          each camera photo the object, the first class, and the median
          latency.
  \end{itemize}
  \keypoint{The file gives the scale and the zero point of its output. Read
    them from the output details. Do not write them into the code.}
\end{frame}

% ------------------------------------------------------------------- Part C
\begin{frame}{Part C: export and inspect (45 min)}
  \small
  Work on the laptop, in the notebook \code{export\_inspect.ipynb}.
  \vskip 0.2em
  \renewcommand{\arraystretch}{1.15}
  \begingroup\centering
  \begin{tabular}{@{}L{0.12\textwidth}L{0.66\textwidth}L{0.11\textwidth}@{}}
    \toprule
    \textbf{Section} & \textbf{Work} & \textbf{Time} \\
    \midrule
    0 and 1 & Setup, and the layers of the PyTorch model & 3 min \\
    2 & \textbf{Task 1:} export to ONNX with a static shape & 7 min \\
    3 & \textbf{Task 2:} export with a dynamic batch size, and a comparison
        of the two files & 8 min \\
    4 & Export to LiteRT, and the operators of the file & 5 min \\
    5 & The graph for each optimization level, and the files in Netron
      & 12 min \\
    6 & \textbf{Task 3:} fold a batch normalization by hand & 7 min \\
    7 and 8 & The same result in each runtime, and the files for the
        board & 3 min \\
    \bottomrule
  \end{tabular}\par\endgroup
  \vskip 0.4em
  \begin{itemize}
    \item The cell after a task prints \code{Task N: complete} or
          \code{not complete}.
    \item The last cell prints \code{Tasks complete: 1, 2, 3}.
    \item At the end, copy four files to the board with the \code{scp}
          command of section 8.
  \end{itemize}
\end{frame}

\begin{frame}{Part C: Tasks 1 and 2, the two exports}
  \small
  \renewcommand{\arraystretch}{1.2}
  \begingroup\centering
  \begin{tabular}{@{}L{0.09\textwidth}L{0.25\textwidth}L{0.56\textwidth}@{}}
    \toprule
    \textbf{Task} & \textbf{Function} & \textbf{You write} \\
    \midrule
    1 & \code{export\_static} & The call of \code{torch.onnx.export}, with
        the names and the options of the notebook \\
    2 & \code{export\_dynamic} & The same call with \code{dynamic\_shapes}:
        dimension 0 of the input gets the name \code{batch} \\
    \bottomrule
  \end{tabular}\par\endgroup
  \vskip 0.4em
  \begin{itemize}
    \item The checks: the static file has the input shape 1, 3, 224, 224.
          The dynamic file accepts 3 images. The outputs agree with
          PyTorch.
    \item \textbf{You record:} the layers of the PyTorch model, the nodes
          of the ONNX file, and the layer type that has no node.
  \end{itemize}
  \begin{exerciseblock}
    \textbf{Prediction,} before the cell that compares the two files:
    \begin{itemize}
      \item Is the dynamic file slower for one image?
      \item What does the static file do with 2 images?
      \item Is a batch of 4 faster for each image than 4 runs with one
            image?
    \end{itemize}
  \end{exerciseblock}
\end{frame}

\begin{frame}{Part C: the graphs in Netron}
  \begingroup\centering
  \begin{tikzpicture}[font=\scriptsize, >={Stealth[length=4pt]},
      g/.style={draw=coolgray, rounded corners=2pt, fill=boxgray,
                minimum width=1.9cm, minimum height=0.6cm, align=center},
      q/.style={draw=cardinalred, rounded corners=2pt, fill=cardinalred!10,
                minimum width=1.9cm, minimum height=0.6cm, align=center}]
    \node[coolgray, anchor=east] at (-1.2,0)    {the PyTorch model};
    \node[g] (a) at (0,0)   {Conv2d};
    \node[g] (b) at (2.4,0) {BatchNorm2d};
    \node[g] (c) at (5.1,0) {ReLU6};
    \draw[->] (a) -- (b); \draw[->] (b) -- (c);
    \node[coolgray, anchor=east] at (-1.2,-0.9) {in each file};
    \node[q] (d) at (0,-0.9)   {?};
    \node[q] (e) at (2.4,-0.9) {?};
    \node[q] (f) at (5.1,-0.9) {?};
    \draw[->] (d) -- (e); \draw[->] (e) -- (f);
  \end{tikzpicture}\par\endgroup
  \vskip 0.4em
  \small
  \renewcommand{\arraystretch}{1.15}
  \begingroup\centering
  \begin{tabular}{@{}L{0.36\textwidth}L{0.56\textwidth}@{}}
    \toprule
    \textbf{File in \code{models/}} & \textbf{Question of the report} \\
    \midrule
    \code{mnv2\_level\_basic.onnx} & Which node is gone? Which new input
                                     does the first \code{Conv} have? \\
    \code{mnv2\_level\_extended.onnx} & What is the name of the first node?
                                     Which attribute holds the activation? \\
    \code{mnv2.tflite} & Which property of the first \code{Conv2D} holds the
                         activation? Which operator comes before it, and
                         why? \\
    \bottomrule
  \end{tabular}\par\endgroup
  \vskip 0.4em
  \begin{itemize}
    \item Drag a file into the page \code{netron.app}. Click a node to see
          its inputs and its attributes.
    \item \textbf{You record:} the table of the four levels with the nodes
          and the latency of your laptop. Write your prediction first:
          which level gives the largest gain?
  \end{itemize}
\end{frame}

\begin{frame}{Part C: Task 3, fold a batch normalization}
  \[
    s = \frac{\gamma}{\sqrt{\sigma^2 + \varepsilon}} \qquad
    W' = s\,W \qquad b' = \beta - s\,\mu
  \]
  \small
  \begin{itemize}
    \item The function \code{fold\_batch\_norm} gets the weight of the first
          convolution and the four tensors of its batch normalization. It
          returns $W'$ and $b'$.
    \item The convolution has no bias, so $b = 0$ in the formula of the
          lecture.
    \item Each output channel has its own $s$. The weight has the shape
          (output channels, input channels, height, width). Use
          \code{s[:, None, None, None]} to multiply each filter by its
          value.
    \item The check runs the two operators of PyTorch, and then one
          convolution with your $W'$ and $b'$. The task is complete when
          the largest difference is below 0.0001.
    \item \textbf{You record:} the largest difference, and the number of
          values that the model loses and gets.
  \end{itemize}
  \keypoint{This is the exercise of the lecture, for 32 filters in place of
    one. The basic level of ONNX Runtime did it for all 52 layers.}
\end{frame}

% ------------------------------------------------------------------- Part D
\begin{frame}[fragile]{Part D: measure (35 min)}
  \small
  \renewcommand{\arraystretch}{1.15}
  \begingroup\centering
  \begin{tabular}{@{}L{0.07\textwidth}L{0.70\textwidth}L{0.12\textwidth}@{}}
    \toprule
    \textbf{Step} & \textbf{Work} & \textbf{Time} \\
    \midrule
    1 & Predict: write your answer for each question of the report
      & 5 min \\
    2 & Task D1: the function \code{median\_latency\_ms} in
        \code{pi/bench.py} & 5 min \\
    3 & Measure: run the script two times & 10 min \\
    4 & Fill the latency table, and calculate the three ratios & 10 min \\
    5 & Compare with your predictions, and write the Decision Log
      & 5 min \\
    \bottomrule
  \end{tabular}\par\endgroup
\begin{lstlisting}
runtime      file             type     threads   load ms  median ms
litert       mnv2.tflite      float32        1       ...        ...
...
onnxruntime  mnv2_int8.onnx   int8           4       ...        ...
\end{lstlisting}
  \begin{itemize}
    \item \code{python pi/bench.py} measures 6 files with 1, 2, and 4
          threads: 18 lines. It also writes \code{results.csv}.
    \item Task D1: some calls to warm up, then one time for each call, then
          the median.
  \end{itemize}
\end{frame}

\begin{frame}{Part D: predict, measure, compare}
  \begin{exerciseblock}
    \textbf{Prediction,} before you run the script:
    \begin{itemize}
      \item \code{float32}: which runtime is faster, LiteRT or ONNX
            Runtime?
      \item The gain from 1 thread to 4 threads: 4, more than 3, or less
            than 3?
      \item LiteRT: is \code{int8} faster than \code{float32}? By which
            factor?
      \item ONNX Runtime: is the dynamic file slower than the static file?
    \end{itemize}
  \end{exerciseblock}
  \small
  \begin{itemize}
    \item Close each other program, and keep \code{htop} open. Use the
          second run for the table.
    \item \textbf{You record:} the latency table, the temperature, the busy
          cores in \code{htop}, and three ratios: threads, precision, and
          runtime.
    \item If you have time: \code{python pi/bench.py --levels} measures
          the four optimization levels of ONNX Runtime on the board.
  \end{itemize}
  \keypoint{A latency is valid for one model, one runtime, one setting, and
    one device. Write all four next to each number.}
\end{frame}
